\documentclass[10pt,a4paper]{amsart}
\usepackage[margin=19mm]{geometry}
\usepackage{amsmath,amssymb,mathtools}
\usepackage[scr=rsfs,cal=euler]{mathalfa}
\usepackage{xfrac}
\usepackage[unicode,hidelinks,bookmarks=false]{hyperref}
\newcommand{\ie}{i.e.\ }
\newcommand{\cf}{cf.\ }
\newcommand{\resp}{respectively\ }
\newcommand{\Subsection}{\subsection}
\providecommand{\nobreakdash}{\nobreak\hbox{-}}
\setlength{\emergencystretch}{2em}
\multlinegap0pt
\usepackage{bm}
\usepackage{bbm}
\usepackage{dsfont}
\usepackage{xspace}
\usepackage{cancel}
\usepackage{mathtools}
\usepackage{amsthm}
\makeatletter
\@ifundefined{thm}{\newtheorem{thm}{Thm}}{}
\@ifundefined{lemma}{\newtheorem{lemma}[thm]{Lemma}}{}
\@ifundefined{prop}{\newtheorem{prop}[thm]{Proposition}}{}
\makeatother
\makeatletter
\newcommand{\eps}{{\epsilon}}
\def\ge{\geqslant}
\expandafter\ifx\csname pfof\endcsname\relax
\newenvironment{pfof}[1]{\vspace{1ex}\noindent{\bf Proof of
#1}\hspace{0.5em}}{\hfill\qed\vspace{1ex}}
\fi
\def\phi{\varphi}
\makeatother
\title[H\"older continuity of measures for heavy tail potentials]{H\"older continuity of measures for heavy tail potentials}
\author{Godofredo Iommi, Dalia Terhesiu, Mike Todd}
\date{Source: arXiv:2308.14485v3 (CC BY 4.0)}
\begin{document}
\maketitle

\noindent\emph{Source and licence.} H\"older continuity of measures for heavy tail potentials. Version arXiv:2308.14485v3 (CC BY 4.0), distributed under the Creative Commons Attribution 4.0 International licence, as recorded in the arXiv licence metadata for this exact version and shown on the abstract page https://arxiv.org/abs/2308.14485v3. Full rights notice: licenses/arxiv-2308.14485-CC-BY-4.0.txt.\par\smallskip

\noindent\textit{Editorial scope.} Counted unit: the Lemma whose source label is \texttt{lemma:doubderder} in the article (source0.tex lines 1535--1554, source label \texttt{lemma:doubderder}), delivered together with the article's own complete original proof of it (source0.tex lines 1555--1607, one \texttt{proof} environment of 53 lines). No mathematics is written, repaired or re-derived: statement and proof are the article's own text line for line. Required context retained uncounted: prop:gmsec (lines 511--518); eq:dounderpbar (lines 949--953); lemm:refev (lines 1185--1253). Every definition, display and label of those lines is retained unchanged. Omitted from this package and not redistributed: 1--510, 519--948, 954--1184, 1254--1534, 1608--2200. Nothing inside a retained excerpt is omitted or altered: every retained body is delivered exactly as the article's lines are written, so no omission receipt of any kind accompanies this package. Citations: the retained excerpts carry no citation command at all, so no bibliography entry of the article is transcribed and no reference is redistributed. Reference adaptation: none was needed and none was made; every reference inside the delivered unit resolves inside it, so no unresolved reference or citation is produced. Printed slips kept literal: no printed word, symbol, hypothesis or number of the counted statement or of its proof was altered. Layout and class: the article's own class is \texttt{article} with packages including amsmath, amssymb, theorem, enumerate, hyperref, bm; the wrapper is the lane's pinned \texttt{amsart} editorial wrapper at 10pt on A4, which restates the theorem-like environments the retained text needs and adds the article's own macro definitions that the retained text needs (\texttt{\textbackslash eps}, \texttt{\textbackslash ge}, \texttt{\textbackslash phi}), with extra wrapper packages (bm, bbm, dsfont, xspace, cancel, mathtools). The delivered numbering is the unit's own. Assets: no retained span contains a figure environment, a graphics-inclusion command, a PDF-page inclusion, a table or a TikZ picture; no image file is copied into this package, the package declares no asset dependency and no image file is redistributed with it. Licence: version arXiv:2308.14485v3 of this article is distributed under the Creative Commons Attribution 4.0 International licence, as recorded in the arXiv licence metadata for this exact version and shown on the abstract page https://arxiv.org/abs/2308.14485v3. A full rights notice is delivered at licenses/arxiv-2308.14485-CC-BY-4.0.txt. Characters: every retained excerpt is pure ASCII, so the delivered engine, XeLaTeX with its default Latin Modern text font, reports no missing character.
\begin{prop}\label{prop:gmsec} Assume (GM0) with $\mu_Y(\tau\ge x)=c x^{-\beta}(1+o(1))$ for $\beta\in (1,2)$.
Suppose that (GM1) holds with $\psi_0=C_1 \tau^\gamma$ with $\gamma\in (\beta-1,1)$. 
There exist $C_2, C_3>0$ depending only on $c,\beta,\gamma$ and $\tau^*$ so that the following hold as $s\to 0$.
\begin{itemize}
\item[(i)] If $\beta/\gamma\in (1,2]$, then $p''(s)=C_2 s^{\beta-\gamma-1}(1+o(1))$.
\item[(ii)] If $\beta/\gamma\in (2, 3)$, then $p'''(s)=-C_3 s^{\beta-2\gamma-1}(1+o(1))$.
\end{itemize}
\end{prop}

 \begin{align}\label{eq:dounderpbar}
  \frac{\partial}{\partial s} D(u,s)&=\frac{\partial}{\partial s}\frac{\partial}{\partial u}\bar p(u,s)
  =  -\frac{\frac{\partial}{\partial u}\lambda(u,s)\frac{\partial}{\partial s}\lambda(u,s)}{\lambda(u,s)^2}+\frac{\frac{\partial}{\partial s}\frac{\partial}{\partial u}\lambda(u,s)}{\lambda(u,s)}\\
 \nonumber &=-\bar\psi^*\tau^*-E_0(u,s)+\frac{\frac{\partial}{\partial s}\frac{\partial}{\partial u}\lambda(u,s)}{\lambda(u,s)},
 \end{align}

\begin{lemma}\label{lemm:refev}
The following hold in the setup of Proposition~\ref{prop:gmsec}.  Let $u, s\in [0,\delta_0)$.
\begin{itemize}
 \item[(i)] $\frac{\partial}{\partial u}\lambda(u,s)=-\tau^*+d(u,s)$, where $d(u,s)$ is as follows.
 
 There exists $C>0$ depending only on $c,\beta$ so that $d(u,0)=
 Cu^{\beta-1}(1+o(1))$.  
 Moreover, there exist $C_2, C_3>0$ depending only on $c,\beta, \gamma$ so that
 as $u,s\to 0$,
 \begin{align*}
 \frac{\partial}{\partial s} d(u,s)=
 C_2 u^{\beta-\gamma-1}(1+o(1)), &\text{ if }\beta/\gamma\in (1,2],
   \end{align*}
 
  \begin{align*}
 \frac{\partial^2}{\partial s^2} d(u,s)=
 C_3 u^{\beta-\gamma-2}(1+o(1)), &\text{ if }\beta/\gamma\in (2,3).
  \end{align*}
  
 
 \item[(ii)] The following holds for some $C, C'>0$ depending only on $c,\beta/\gamma$.
  \begin{align*}
\frac{\partial}{\partial s}\lambda(u,s)=\bar\psi^*+ e(u,s) + \begin{cases}
h(s)+ h_0(s), &\text{ if }\beta/\gamma\in (1,2],\\
 -s\int_Y \bar\psi^2\, d\mu_{\overline{\phi}}+ Cs^{\beta/\gamma-1}+h_1(s), &\text{ if }\beta/\gamma\in (2,3),
 \end{cases}
  \end{align*} 
  where $h(s)=Cs^{\beta/\gamma-1}$ if  $\beta/\gamma\in (1,2)$, $h(s)=C\log(1/s)$ if  $\beta/\gamma=2$
  and where
    $h_0$, $h_1$  and $e$ are as follows: 
 \begin{itemize}
 \item[(a)] $h_0(s)=o(s^{\beta/\gamma-1})$, $h'_0(s)=o(s^{\beta/\gamma-2})$ if $\beta/\gamma\in (1,2)$
 and $h'_0(s)=o(\log(1/s))$ if $\beta/\gamma=2$.
 
 \item[(b)] $h_1(s)=o(s^{\beta/\gamma-1})$, $h'_1(s)=C' s^{\beta/\gamma-2}(1+o(1))$
  and $h''_1(s)=C' s^{\beta/\gamma-3}(1+o(1))$;
  
 \item[(c)] $e(0,s)=O(s)$, $e(u,0)=o(1)$ as $u,s\to 0$ and
  
 
  \begin{itemize}
  \item[(*)] If $\beta/\gamma\in (1,2)$, then $\frac{\partial}{\partial s} e(u,s)= o(u^{\beta-\gamma-1})+o(s^{\beta/\gamma-2})$.
  Also, $\frac{\partial}{\partial s} e(0,s)=o(s^{\beta/\gamma-2})$.
   
   \item[(**)] If $\beta/\gamma=2$, then $\frac{\partial}{\partial s} e(u,s)= o(u^{\beta-\gamma-1})+o(\log(1/s))$. Also, $\frac{\partial}{\partial s} e(0,s)=o(\log(1/s))$.
    
   \item[(***)] If $\beta/\gamma\in (2,3)$, then $\frac{\partial}{\partial s} e(u,s)= o(u^{\beta-\gamma-2})+o(s^{\beta/\gamma-3})$.   Also, $\frac{\partial}{\partial s} e(0,s)= o(s^{\beta/\gamma-3})$.    \end{itemize}
    \end{itemize}

  
  \item[(iii)] Let
  $\kappa(u,s)=\frac{\partial}{\partial s}\frac{\partial}{\partial u}\lambda(u,s)$.
  Then there exist $C, C'>0$ depending only on $c,\beta,\gamma$,
  so that
  $$
  \begin{cases}
  \kappa(u,0)=C u^{\beta-\gamma-1} + O(u^{\beta-\gamma-1+\eps_0}), & \text{ if } \beta/\gamma\in (1,2], \\
  \frac{\partial}{\partial s}\kappa(u,s)\Big|_{s=0}=C' u^{\beta-2\gamma-1} + O(u^{\beta-2\gamma-1+\eps_0}), & \text{  if } \beta/\gamma\in (2,3),
  \end{cases}
  $$
   as $u\to 0$ and for any $\eps_0>0$.
   
   
   Also, the following hold for some $\hat C_2, \hat C_3>0$ depending only on $c,\beta,\gamma$, as $u,s\to 0$.
   \begin{itemize}
    \item[(*)] If $\beta/\gamma\in (1,2]$, then $\frac{\partial}{\partial s}\kappa(u,s)=\hat C_2 u^{\beta-2\gamma-1}(1+o(1))$.
    \item[(**)] If $\beta/\gamma\in (2,3)$, then $\frac{\partial^2}{\partial s^2}\kappa(u,s)=-\hat C_3 u^{\beta-3\gamma-1}(1+o(1))$.    \end{itemize}
    \end{itemize}
\end{lemma}

\begin{lemma}\label{lemma:doubderder} Assume the setup of Proposition~\ref{prop:gmsec} with larger range of $\gamma$,
namely $\gamma\in (\beta-1,\beta)$. There exist $C_2, C_3, C_4, C_5>0$(varying from line to line) so that the following hold as $u,s\to 0$.
\begin{itemize}
\item[(i)] If $\beta/\gamma\in (1,2)$ then 
\(
\frac{\partial^2}{\partial s^2}\frac{\partial}{\partial u}\bar p(u,s)=-C_2s^{\beta/\gamma-2}(1+o(1)+C_3 u^{\beta-2\gamma-1}(1+o(1)).\)
Also, $\frac{\partial}{\partial s}\frac{\partial}{\partial u}\bar p(u,s)= C_4u^{\beta-\gamma-1}(1+o(1))+C_5 s u^{\beta-2\gamma-1}(1+o(1))$.


\item[(ii)] If $\beta/\gamma=2$ then 
\(
\frac{\partial^2}{\partial s^2}\frac{\partial}{\partial u}\bar p(u,s)=-C_2\log(1/s)(1+o(1))+C_3 u^{-1}(1+o(1)).\)
Also, $\frac{\partial}{\partial s}\frac{\partial}{\partial u}\bar p(u,s)= C_4u^{\beta-\gamma-1}(1+o(1))+ C_3s u^{-1}(1+o(1))-C_2 s\log(1/s)(1+o(1))$.
\item[(ii)] If $\beta/\gamma\in (2,3)$ then 
\(
\frac{\partial^3}{\partial s^3}\frac{\partial}{\partial u}\bar p(u,s)=-C_2s^{\beta/\gamma-3}(1+o(1))-C_3 u^{\beta-3\gamma-1}(1+o(1)).\)
Also, $\frac{\partial^2}{\partial s^2}\frac{\partial}{\partial u}\bar p(u,s)\Big|_{s=0}=- C_4u^{\beta-2\gamma-1}(1+o(1))+C_3 s u^{\beta-3\gamma-1}(1+o(1))$.
\end{itemize}

\end{lemma}

\begin{proof} First we recall from~\eqref{eq:dounderpbar} that
\begin{align*}
 \frac{\partial}{\partial s}\frac{\partial}{\partial u}\bar p(u,s)
  &=  -\frac{\frac{\partial}{\partial u}\lambda(u,s)\frac{\partial}{\partial s}\lambda(u,s)}{\lambda(u,s)^2}+\frac{\frac{\partial}{\partial s}\frac{\partial}{\partial u}\lambda(u,s)}{\lambda(u,s)}.
\end{align*}
Set $A(u,s):=\frac{\partial}{\partial u}\lambda(u,s)\frac{\partial}{\partial s}\lambda(u,s)$ and recall (for instance, from Lemma~\ref{lemm:refev}(iii)) that 
$\kappa(u,s)=\frac{\partial}{\partial s}\frac{\partial}{\partial u}\lambda(u,s)$.
Compute that
\begin{align*}
 \frac{\partial^2}{\partial s^2}\frac{\partial}{\partial u}\bar p(u,s)
  &=  -\frac{\frac{\partial}{\partial s}A(u,s)}{\lambda(u,s)^2}-2\frac{A(u,s)\frac{\partial}{\partial s}\lambda(u,s) }{\lambda(u,s)^3}
  +\frac{\frac{\partial}{\partial s}\kappa(u,s)}{\lambda(u,s)}-\frac{\kappa(u,s)}{\lambda(u,s)^2}\\
  &=:N_1(u,s)+N_2(u,s)+N_3(u,s)+N_4(u,s).
\end{align*}


We provide the proof of item (i). Item (ii) follows by the same argument using the statements for the case $\beta/\gamma=2$ in Lemma~\ref{lemm:refev} (i) and (ii). Item (iii) follows by a similar argument after differentiating once more and using the statements for the case $\beta/\gamma\in (2,3)$ in Lemma~\ref{lemm:refev} (i) and (ii). 

From the estimates of Lemma~\ref{lemm:refev} (i) and (ii) (the statements for the case $\beta/\gamma\in (1,2)$), it is easy to see that $N_2$ and $N_4$
do not contribute to the main asymptotics (because they go to a constant as $u,s\to 0$). We need to look at  $N_1$ and $N_3$.\\


 
 
\textbf{The term $N_1(u,s)$}.
Using the same notation as in Lemma~\ref{lemm:refev} (i) and (ii),
\begin{align*}
A(u,s)=
\left(-\tau^*+d(u,s)\right)\left(\bar\psi^*+ C s^{\beta/\gamma-1}+h(s)+h_0(s) +e(u,s) \right)
\end{align*}
and 
\begin{align*}
\frac{\partial}{\partial s} A(u,s)=&
\frac{\partial}{\partial s} d(u,s)\, \left(\bar\psi^*+ C s^{\beta/\gamma-1}+h(s)+h_0(s)  +e(u,s) \right)\\
&+ \left(-\tau^*+d(u,s)\right)\left( C(\beta/\gamma-1) s^{\beta/\gamma-2}+\frac{\partial}{\partial s} h(s)+h_0'(s)  +\frac{\partial}{\partial s} e(u,s) \right).\end{align*}


Using all the estimates on $d, h_0, e$ in Lemma~\ref{lemm:refev} (i) and (ii) (the statements for the case $\beta/\gamma\in (1,2)$),
we obtain that there exist $C_2, C_2'>0$ so that
\begin{align*}
\frac{\partial}{\partial s} A(u,s)&=-C_2 s^{\beta/\gamma-2}(1+o(1)) + C_2' u^{\beta-\gamma-1}(1+o(1)),
\end{align*}
which gives the asymptotics for $N_1(u,s)$.
In the previous displayed equation, apart from the estimates on $\frac{\partial}{\partial s} d(u,s)$ and  $\frac{\partial}{\partial s} e(u,s)$, we have used the immediate consequence of 
 Lemma~\ref{lemm:refev}(ii) that $d(u,s)=O(s u^{\beta-\gamma-1})$ and that $ e(u,s)= o(s u^{\beta-\gamma-1})$.
 
 \textbf{The term $N_3(u,s)$}. By Lemma~\ref{lemm:refev} (iii) (the statement for the case $\beta/\gamma\in (1,2)$),
 $\frac{\partial}{\partial s}\kappa(u,s)=C_3 u^{\beta-2\gamma-1}(1+o(1))$, for some $C_3>0$. This gives the same asymptotics for $N_3$. Therefore,
\begin{align*}
N_1(u,s)+N_3(u,s)=-C_2 s^{\beta/\gamma-2}(1+o(1))+C_3 u^{\beta-2\gamma-1}(1+o(1)),\end{align*}
which gives the first statement in item (i).

The second statement in item (i) follows immediately from the first together with the asymptotics of $\kappa(u,0)$ in Lemma~\ref{lemm:refev} (iii).~\end{proof}

\end{document}
